% !TEX root = ../Main.tex
\chapter{极限思想}

\section{从平均速率到瞬时速率}

平均速率就是总路程除以总时间：
\[
\bar v=\frac{S}{t}\qquad\Big(\text{更一般地，}\ \bar v=\frac{\Delta S}{\Delta t}\Big).
\]

在位移－时间（$S$–$t$）图上，取 $A$、$B$ 两点，$\bar v=\dfrac{\Delta S}{\Delta t}$ 正是\textbf{割线} $AB$ 的斜率 $k_{AB}$。

\begin{center}
\begin{tikzpicture}[>=Stealth, scale=1.1]
    \draw[->] (0,0) -- (0,3) node[above] {$S$};
    \draw[->] (0,0) -- (4,0) node[right] {$t$};
    \draw[thick, blue, domain=0:2.6, samples=60, smooth] plot (\x, {0.42*\x*\x});
    \coordinate (A) at (1.3,{0.42*1.3*1.3});
    \coordinate (B) at (2.4,{0.42*2.4*2.4});
    \filldraw (A) circle (1.2pt) node[below right=1pt] {$A$};
    \filldraw (B) circle (1.2pt) node[above right=1pt] {$B$};
    \draw[dashed] (B) -- (1.3,{0.42*2.4*2.4}) node[midway, above] {\footnotesize $\Delta S$};
    \draw[dashed] (1.3,{0.42*2.4*2.4}) -- (A);
    \node[below=2pt] at (1.85,{0.42*1.3*1.3}) {\footnotesize $\Delta t$};
    \draw[red, thick] ($(A)!-0.3!(B)$) -- ($(A)!1.3!(B)$);
\end{tikzpicture}
\end{center}

现在\textbf{减小 $\Delta t$}：让 $B$ 沿曲线滑向 $A$。割线斜率 $k_{AB}$ 随之变化，
\[
\Delta t\downarrow\ \Longrightarrow\ k_{AB}\to k_{\text{切}@A}.
\]
当 $\Delta t\to 0$ 时，$\bar v$ 就趋于 $A$ 点的\textbf{瞬时速率} $v_t$。反映在图上：瞬时速率就是该点处\textbf{切线的斜率}。

\state{极限思想的结论}{
    位移－时间图上，某点切线的斜率 $=$ 该点的瞬时速率。把 $\Delta t$ 不断缩小、让割线逼近切线，就是"极限"。
}

\section{补充：切线为什么用斜率定义}

曲线某点切线的"倾斜程度"，描述的正是曲线在该点的\textbf{变化率}。这件事可以用\textbf{反证法}佐证：

\sol{
    若两者不相等——即所谓的"切线"在该点的斜率与曲线在该点的变化率不一致——那么这条直线在该点附近就会与曲线相交于另一个点，于是它不再是"只碰一下"的切线，矛盾（图中打 $\times$）。所以二者必相等：切线斜率就是曲线在该点的变化率。
}

\begin{center}
\begin{tikzpicture}[>=Stealth, scale=1]
    \draw[thick, blue, domain=-1.5:1.5, samples=60, smooth] plot (\x, {0.7*\x*\x - 0.4});
    \coordinate (P) at (0,-0.4);
    \filldraw (P) circle (1.2pt);
    \node[below left] at (P) {$B$};
    % 错误的"切线"——斜的直线会再次穿过曲线
    \draw[red] (-1.5,{0.7*(-1.5)*(-1.5)-0.4-0.9}) -- (1.5,{0.7*1.5*1.5-0.4+0.9});
    \node[red] at (1.0,0.9) {\footnotesize $\times$};
    % 正确切线——水平
    \draw[thick] (-1.6,-0.4) -- (1.6,-0.4) node[right] {\footnotesize 切线};
\end{tikzpicture}
\end{center}
